% =====================================================================
\section{Итоговый проект: от идеи до защиты}
% =====================================================================
\lessonintro
  {довести собственную идею до работающего устройства и уметь рассказать о нём за пять минут}
  {все предыдущие уроки; опирайтесь на те, что нужны вашей теме}
  {готовое устройство, электрическая схема, описание проекта и выступление на защите}
  {детали выбранной темы (приложение B), компьютер для документации, 4--6 недель}

\subsection{Сначала: проект --- это история с началом и концом}

В уроках вы делали отдельные опыты, каждый из которых отвечал на один вопрос. Проект --- это когда
несколько опытов складываются в устройство, у которого есть задача, пользователь и результат. В нём
мало просто «что-то мигает»: нужно понять, \emph{что} устройство решает, \emph{как} оно это делает и
\emph{как проверить}, что получилось. Такой урок не про новый код, а про умение довести дело до конца.

\subsection{Как выбрать тему}

Хорошая тема объединяет три условия: она вам интересна, решаема за 4--6 недель и опирается на уроки,
которые вы уже освоили. Ниже --- идеи на выбор. Можно взять свою, но проверьте её тем же способом: по
таблице подсчитайте, на какие уроки она опирается и какие новые детали нужны.

\begin{center}\small
\begin{tabularx}{\textwidth}{@{}r X X l@{}}
\toprule
№ & Идея & Что делает устройство & Уроки \\
\midrule
1  & Метеостанция & измеряет температуру и влажность, показывает на OLED и пишет журнал на SD & 5, 6, 10, 15 \\
2  & Парковочный радар & измеряет расстояние до препятствия, показывает лентой и звуком & 4, 15, 17 \\
3  & Умная теплица & включает полив по влажности и освещению, ведёт журнал & 5, 9, 15, 16 \\
4  & Сигнализация & реагирует на движение или дверь и присылает сообщение по MQTT & 3, 8, 12, 15 \\
5  & Пульт для презентаций & листает слайды кнопками и регулирует громкость & 3, 5, 19 \\
6  & Лампа-настроение & лента меняет цвета эффектами, а режим выбирается с телефона & 8, 11, 17 \\
7  & Сетевые часы & показывают московское время на OLED и ленте, яркость по освещённости & 5, 6, 17, 18 \\
8  & Робот-тележка & едет и объезжает препятствия по дальномеру & 4, 15, 16 \\
9  & Камера наблюдения за растением & снимает кадры раз в час и показывает поток по запросу (с согласия окружающих) & 8, 10, 18, 20 \\
10 & Автономный датчик & спит от батарейки, просыпается и отправляет данные другой плате & 9, 12, 15 \\
11 & Умный вентилятор & мотор включается по температуре, порог хранится в памяти & 4, 11, 15, 16 \\
12 & Игра «Реакция» & лента зажигается случайно, нужно успеть нажать кнопку; рекорд сохраняется & 3, 11, 17 \\
\bottomrule
\end{tabularx}
\end{center}

\subsection{Проекты на возможностях ESP32-S3}
После расширенного модуля можно выбрать аудиорегистратор I²S (29), измеритель импульсов RMT/PCNT
(30--31), сеть датчиков CAN/RS-485 (32), цветную панель (33), USB-пульт с консолью CDC (34),
датчик с ULP-пробуждением (36) или распознавание жестов IMU (37).
Для каждой темы задайте измеримый результат: долю потерянных блоков аудио, погрешность частоты,
время восстановления связи, задержку кадра, расход энергии или матрицу ошибок модели.
Испытания из урока 38 включите в защиту. Для локальной панели используйте ESPUI (39),
для выдачи настроек сети --- отдельный этап WiFiManager (40). Не требуется использовать всю периферию одновременно:
правильно обоснованный выбор и подтверждённый результат важнее количества модулей.

\subsection{Техническое задание на полстраницы}

ТЗ --- это договор с самим собой. Оно отвечает на вопрос «что считается готовым проектом». Напишите его
до покупки деталей и первой строчки кода. Используйте такую форму:

\begin{enumerate}
  \item \textbf{Название} и одно предложение: для кого и зачем это устройство.
  \item \textbf{Функции}: 3--5 пунктов вида «устройство должно\ldots».
  \item \textbf{Входы и выходы}: какие датчики и кнопки измеряют мир, какие исполнители им отвечают.
  \item \textbf{Условия и ограничения}: питание, размеры, безопасность (напряжение, ток, съёмка людей).
  \item \textbf{Критерий успеха}: измеримая проверка, например «температура отличается от термометра не больше чем на $1\,$°C».
  \item \textbf{Состав}: список деталей (используйте приложение B).
\end{enumerate}

\subsection{Блок-схема устройства}

Блок-схема помогает увидеть проект целиком до сборки. Рисуйте слева входы, справа выходы, в центре ESP32-S3,
а питание и связь --- отдельно. Для каждой стрелки подпишите сигнал и вывод платы.

\begin{figure}[H]
\centering
\begin{tikzpicture}[font=\sffamily\scriptsize,
  box/.style={blk,minimum width=3.3cm,minimum height=1cm}]
  \node[box,fill=blkper] (in1) at (0,1.5) {Датчики\\ температура, расстояние};
  \node[box,fill=blkper] (in2) at (0,0) {Кнопки,\\ потенциометр};
  \node[box,fill=blkcpu,minimum width=3.6cm,minimum height=3.6cm] (mc) at (6,0.4) {\textbf{ESP32-S3}\\ программа,\\ логика проекта};
  \node[box,fill=blkper] (out1) at (12,1.5) {Исполнители\\ мотор, реле, лента};
  \node[box,fill=blkper] (out2) at (12,0) {Дисплей,\\ звук};
  \node[box,fill=blkrf] (net) at (6,3.4) {Связь\\ Wi-Fi, BLE, MQTT};
  \node[box,fill=blkmem] (mem) at (12,3.4) {Память\\ NVS, SD, LittleFS};
  \node[box,fill=blkrtc] (pw) at (6,-2.5) {Питание\\ USB, батарея, блок \SI{5}{V}};
  \draw[arr] (in1.east) -- (mc.west |- in1.east);
  \draw[arr] (in2.east) -- (mc.west |- in2.east);
  \draw[arr] (mc.east |- out1.west) -- (out1.west);
  \draw[arr] (mc.east |- out2.west) -- (out2.west);
  \draw[darr] (mc.north) -- (net.south);
  \draw[darr] (net.east) -- (mem.west);
  \draw[arr,thick] (pw.north) -- (mc.south);
\end{tikzpicture}
\caption{Шаблон блок-схемы: входы слева, выходы справа, связь и память сверху, питание снизу.
Подпишите на стрелках сигналы и номера GPIO вашего устройства}
\end{figure}

\subsection{План по неделям}

Рассчитывайте на пять-шесть недель. Самая частая причина провала --- пустая последняя неделя,
в которую уходит «всё остальное».

\begin{figure}[H]
\centering
\begin{tikzpicture}[x=1.2cm,y=0.75cm,font=\sffamily\scriptsize]
  \foreach \w in {1,...,6}
    \node[font=\sffamily\scriptsize\bfseries] at ({\w*1.2+3.4},0.4) {нед. \w};
  \foreach \y/\t/\s/\e/\c in {
    -1/{Идея и ТЗ, закупка деталей}/1/1/accent,
    -2/{Схема, каждый узел отдельным скетчем}/2/2/okgreen,
    -3/{Основная программа}/3/3/esp,
    -4/{Сборка, корпус, интеграция}/4/4/warn,
    -5/{Тестирование, документация}/5/5/accent,
    -6/{Репетиция и защита}/6/6/okgreen}{
    \node[anchor=east] at (3.5,\y) {\t};
    \fill[\c!70] ({\s*1.2+2.8},{\y-0.3}) rectangle ({\e*1.2+4.0},{\y+0.3});
  }
  \draw[gray!50] (3.6,0) -- (3.6,-6.5);
  \draw[gray!50] (11.2,0) -- (11.2,-6.5);
\end{tikzpicture}
\caption{План проекта на шесть недель. К пятой неделе устройство должно уже работать:
остаётся время на проверку и документацию}
\end{figure}

\subsection{Чек-лист готовности}

Перед защитой пройдите по списку. Если хоть один пункт пустой, это тема для разговора с педагогом.

\begin{itemize}[leftmargin=*,label=\ding{111}]
  \item Устройство включается от своего источника питания и работает без компьютера.
  \item После выключения и повторного включения оно возвращается к рабочему состоянию.
  \item Каждая функция из ТЗ проверена, результат записан.
  \item Электрическая схема нарисована и совпадает с реальной сборкой.
  \item Код разложен по функциям, в нём нет паролей, а комментарии объясняют, \emph{почему}.
  \item Мощные нагрузки подключены через ключ или драйвер, питание продумано (урок 16).
  \item Подготовлены запасной вариант демонстрации (видео или фото) и копия проекта на втором носителе.
  \item Есть описание проекта на 1--2 страницы: цель, схема, инструкция, выводы.
\end{itemize}

\subsection{Защита за пять минут}

Пять минут --- это мало, поэтому структура строгая. Говорите о главном и не читайте текст с листа.

\begin{figure}[H]
\centering
\begin{tikzpicture}[x=2.8cm,font=\sffamily\scriptsize]
  \foreach \s/\e/\t/\c in {0/0.5/{Зачем}/accent,0.5/1.5/{Как устроено}/okgreen,1.5/3/{Демонстрация}/esp,3/4/{Трудности}/warn,4/5/{Итоги}/accent}{
    \pgfmathsetmacro{\w}{(\e-\s)*2.8-0.08}
    \node[blk,fill=\c!15,draw=\c,minimum height=1.1cm,minimum width=\w cm,inner sep=1pt,align=center,font=\sffamily\scriptsize] at ({(\s+\e)/2},0) {\t};
  }
  \draw[->,>=Stealth] (0,-0.9) -- (5.1,-0.9);
  \foreach \m in {0,1,...,5}
    \draw (\m,-0.8) -- (\m,-1.0) node[below]{\m\ мин};
\end{tikzpicture}
\caption{Распределение времени на защите проекта}
\end{figure}

\begin{center}\small
\begin{tabularx}{\textwidth}{@{}l l X@{}}
\toprule
Время & Блок & Что сказать \\
\midrule
0:00--0:30 & Зачем & какую задачу решает устройство и для кого \\
0:30--1:30 & Как устроено & блок-схема, главные узлы, какие уроки пригодились \\
1:30--3:00 & Демонстрация & устройство работает вживую по сценарию из ТЗ \\
3:00--4:00 & Трудности & одна ошибка, как вы её нашли и исправили \\
4:00--5:00 & Итоги & что получилось, что улучшить, ответы на вопросы \\
\bottomrule
\end{tabularx}
\end{center}

\subsection{Критерии оценки}

\begin{center}\small
\begin{tabularx}{\textwidth}{@{}l c X@{}}
\toprule
Критерий & Баллы & Что оценивается \\
\midrule
Работоспособность & 0--30 & устройство выполняет все функции из ТЗ, стабильно работает без вмешательства \\
Код & 0--20 & понятные имена, функции, комментарии; нет бесконечных ожиданий и паролей в тексте \\
Схема и сборка & 0--15 & схема читаема и совпадает с устройством; питание и земля продуманы \\
Документация & 0--15 & ТЗ, схема, инструкция, результаты проверок, список деталей \\
Выступление & 0--20 & укладывается в 5 минут, понятно рассказывает, отвечает на вопросы \\
\midrule
\textbf{Итого} & \textbf{100} & 85 и более --- отлично; 65--84 --- хорошо; 40--64 --- удовлетворительно \\
\bottomrule
\end{tabularx}
\end{center}

\begin{caution}
\textbf{Типичные ошибки проектов.}
\begin{itemize}[leftmargin=*]
  \item \textbf{Слишком большая идея.} «Умный дом целиком» не получится за пять недель. Выберите одну комнату, одну функцию и доведите её.
  \item \textbf{Всё в последний день.} Сборка и отладка занимают вдвое больше времени, чем кажется. Не оставляйте подключение на финал.
  \item \textbf{Проект без схемы.} «Помню, что куда подключено» перестаёт работать через неделю. Фотографируйте и рисуйте схему сразу.
  \item \textbf{Питание «от чего придётся».} Мотор и лента от порта ноутбука приводят к перезагрузкам. Проверьте ток заранее.
  \item \textbf{Демонстрация без запасного плана.} Wi-Fi в зале может не работать. Подготовьте видео и снимки.
  \item \textbf{Пароли в коде.} Не показывайте и не публикуйте SSID и пароль; перед публикацией замените их заглушками.
\end{itemize}
\end{caution}

\begin{summary}
\begin{itemize}[leftmargin=*]
  \item Хороший проект начинается с ТЗ и измеримого критерия успеха, а не с покупки деталей.
  \item План, схема и чек-лист нужны, чтобы к защите устройство, документация и рассказ были готовы.
  \item Защита --- это пять минут: зачем, как устроено, демонстрация, трудности, итоги.
\end{itemize}
\end{summary}

\begin{selfcheck}
\begin{enumerate}
  \item Какие три условия должна выполнять хорошая тема проекта?
  \item Что такое критерий успеха и чем он отличается от списка функций?
  \item Что должно быть на блок-схеме устройства и почему на ней отдельно рисуют питание?
\end{enumerate}
\end{selfcheck}

\begin{tasks}
\begin{enumerate}
\item \lv{1} Выберите тему из таблицы или предложите свою.
\item \lv{1} Запишите три требования к устройству: что оно должно делать, как это проверить, чем измерить.
\item \lv{2} Напишите ТЗ на полстраницы по форме урока.
\item \lv{2} Нарисуйте блок-схему и электрическую схему устройства с номерами GPIO.
\item \lv{3} Составьте перечень деталей и план по неделям с контрольными точками.
\item \lv{4} Проведите репетицию защиты на 5 минут с таймером, запишите замечания слушателей и исправьте слабые места.
\end{enumerate}
\end{tasks}
